\section{Modus operandi}
\label{sec:modus-operandi}

Un primo aspetto del metodo di lavoro ha riguardato il confronto tra la conoscenza di dominio condivisa dai referenti aziendali e le evidenze emerse dai dati. Le indicazioni ricevute hanno guidato la comprensione del processo produttivo e la formulazione di ipotesi, poi verificate sul dataset disponibile. Quando sono emerse discrepanze, abbiamo approfondito l'analisi e chiesto chiarimenti, per comprenderne le cause e precisare le condizioni in cui le ipotesi risultavano valide.

Su questa base, nella preparazione dei dati abbiamo dato precedenza alla coerenza logica dei collegamenti tra le operazioni e dei passaggi di materiale. Abbiamo prima verificato il rispetto dei bilanci di massa e la compatibilità temporale delle lavorazioni.

La documentazione aziendale descrive percorsi che possono diramarsi e ricongiungersi, mentre i dati forniti non identificano esplicitamente i singoli rami e il materiale che li attraversa. Per il modello abbiamo quindi adottato una rappresentazione sequenziale di ciascun ODP, assumendo che il materiale già conforme attenda il completamento del percorso della parte da recuperare prima dell'approntamento dell'ordine. Questa rappresentazione costituisce un'ipotesi di modellazione e non una ricostruzione esplicita dei rami a partire dai consuntivi.

Quando i dati e la conoscenza di dominio lasciavano ulteriori ambiguità, abbiamo introdotto ipotesi di lavoro per ottenere ricostruzioni plausibili, distinguendole dalle indicazioni aziendali e dai fatti osservati. Solo in seguito abbiamo svolto le analisi statistiche, evitando di considerare la sola rarità di un valore come motivo sufficiente per correggerlo.
